\documentclass{situationist-preprint}

\title{Beauty as a Response-Dependent Property}
\msid{AES-02-a}
\msfamily{AES-02}
\msdiscipline{Aesthetics}
\mstype{Principal manuscript}
\mskeywords{beauty; response-dependence; secondary qualities; aesthetic realism; subjectivism; standard of taste; qualified observers}

\begin{document}
\maketitle

\begin{abstract}
The saying that beauty is in the eye of the beholder expresses subjectivism: an object is beautiful for a person just in case it pleases that person, and there is nothing further to be right or wrong about. Its realist opponent holds that beauty is a property objects have independently of anyone's responses. We argue that both views are false and that beauty is a response-dependent property: an object is beautiful just in case it is disposed to give a distinctive pleasure to qualified observers in suitable conditions. We state four conditions any account of beauty must meet: answerability to experience, correctability, supervenience on the non-aesthetic features of the object, and non-circularity. We show that subjectivism fails correctability, that response-independent realism cannot explain answerability, and that a response-dependent account, with its observers and conditions specified on Humean lines, meets all four. It follows that beauty is located neither in the object alone nor in the beholder, but in the relation between the object and a kind of observer. The result is a relational realism about beauty modeled on the corresponding view of colour.
\end{abstract}

\section{The problem}

\enquote{Beauty is in the eye of the beholder} has two readings. On a weak reading it says only that beauty must be experienced to be appreciated, which no one denies. On the strong reading it says that whether an object is beautiful is settled by the response of whoever is looking at it: the rose is beautiful for you if it pleases you, and not beautiful for me if it does not, and neither of us is mistaken. The question we address is whether the strong reading is true. We argue that it is not, and that the truth lies in a position the slogan obscures.

The modern debate begins with \citet{Hume1757}, who held that beauty is no quality in things themselves but exists in the mind that contemplates them, and who nevertheless argued that there is a standard of taste, fixed by the joint verdict of critics who meet certain conditions. \citet{Kant1790} held that the judgment of taste is grounded in a feeling of pleasure, and so is not a cognitive judgment about the object, but that it claims the agreement of everyone, as a judgment of mere agreeableness does not. Both thus combine a dependence of beauty on feeling with a claim to correctness that subjectivism denies. \citet{Santayana1896} gave the subjectivist side its classic formulation, defining beauty as pleasure regarded as the quality of a thing.

Twentieth-century analytic aesthetics reopened the question through the study of aesthetic concepts. \citet{Sibley1959} argued that the application of aesthetic concepts requires taste, a capacity beyond normal perception, and is not governed by conditions stated in non-aesthetic terms, although aesthetic properties depend on non-aesthetic ones. \citet{Pettit1983} argued that aesthetic characterizations can be true or false, while holding that aesthetic properties are, like colours, properties whose instances must be perceptually encountered to be known. In metaethics, \citet{McDowell1985} defended an analogy between values and secondary qualities, on which values are genuine features of the world whose nature is to elicit certain responses; and \citet{Johnston1989} and \citet{Wright1992} developed the general theory of response-dependent concepts. Against these views, \citet{Zangwill2001} has defended a more robust aesthetic realism, on which beauty is a property of objects that our pleasure tracks rather than constitutes. \citet{Goldman1995} and \citet{Levinson2001}, from different directions, have argued that ideal observers may diverge, so that some aesthetic judgments are true only relative to a sensibility.

Three positions are therefore on the table. \emph{Subjectivism} says that beauty is relative to each individual judge. \emph{Response-independent realism} says that beauty is a monadic property of objects, fixed independently of responses. \emph{Response-dependence} says that beauty is a property of objects whose nature consists in a relation to the responses of a kind of observer. We defend the third.

\section{Framework}

We write $B(x)$ for \enquote{$x$ is beautiful}. The pleasure in question is the distinctive pleasure of finding something beautiful, which since Kant has been characterized as disinterested: it is taken in the object as perceived or contemplated, not in its possession, use or consequences. We call it \emph{aesthetic pleasure}.

\begin{definition}[Subjectivism]\label{d:subj}
Subjectivism holds that, for every object $x$ and judge $s$, $x$ is beautiful for $s$ if and only if $x$ gives $s$ aesthetic pleasure, and that there is no non-relative fact about whether $x$ is beautiful.
\end{definition}

\begin{definition}[Response-independent realism]\label{d:real}
Response-independent realism holds that beauty is a property of objects whose possession does not consist in, and is not a priori connected to, any disposition of the object to affect observers.
\end{definition}

\begin{definition}[Response-dependence]\label{d:rd}
A property $F$ is \emph{response-dependent} if there is a response $R$, a kind of observer $O$ and a kind of condition $C$ such that it is a priori that $x$ has $F$ if and only if $x$ is disposed to elicit $R$ in observers of kind $O$ in conditions of kind $C$; and $O$ and $C$ can be specified without using the concept of $F$.
\end{definition}

The final clause of Definition~\ref{d:rd} is what distinguishes a response-dependent property from a trivial one. Any property satisfies a biconditional of the form \enquote{$x$ is $F$ iff $x$ would be judged $F$ by observers who judge correctly}; the biconditional is informative only if correctness is not built into the specification of the observers \citep{Johnston1989,Wright1992}.

\begin{definition}[Qualified observer]\label{d:qual}
A \emph{qualified observer} is one with normal perceptual capacities in the relevant modality, delicacy of discrimination, practice with objects of the relevant kind, acquaintance with a range of comparable objects, and freedom from interests and prejudices bearing on the object. \emph{Suitable conditions} are conditions in which the object is perceived fully and attentively, in its relevant category, without distraction.
\end{definition}

Definition~\ref{d:qual} follows the conditions Hume gave for the true judge: strong sense, delicate sentiment, practice, comparison, and freedom from prejudice \citep{Hume1757}. Each condition is specified without reference to beauty. Delicacy, for example, is the capacity to discriminate fine differences in the non-aesthetic features of an object; practice and comparison are matters of exposure; freedom from prejudice is the absence of interests in the object's owner, author or consequences.

Any account of beauty must respect four conditions.

\begin{principle}[Answerability]\label{p:answer}
It is a priori that a judgment that $x$ is beautiful is answerable to the aesthetic pleasure that $x$ gives when perceived: there is no possibility, knowable a priori to be coherent, of an object that is beautiful and is disposed to give no aesthetic pleasure to any observer however qualified.
\end{principle}

\begin{principle}[Correctability]\label{p:correct}
A sincere judgment that $x$ is beautiful, made by a judge to whom $x$ gives aesthetic pleasure, can be mistaken; and when two judges disagree about whether $x$ is beautiful, at least one of them may be mistaken.
\end{principle}

\begin{principle}[Supervenience]\label{p:super}
Whether an object is beautiful supervenes on its non-aesthetic properties, together with the category in which it is properly perceived.
\end{principle}

\begin{principle}[Non-circularity]\label{p:noncirc}
An account of beauty must not use the concept of beauty in specifying what makes an object beautiful.
\end{principle}

Principle~\ref{p:answer} records the link between beauty and pleasure that runs through Hume, Kant and Sibley, and the related point that aesthetic judgments, unlike judgments of mass, are not properly made on testimony alone. A person who says that a melody is beautiful but that no one, however attentive and experienced, could enjoy hearing it has not described a recondite beauty; she has failed to understand what she said. Principle~\ref{p:correct} records the practice of criticism. We revise our aesthetic judgments in the light of more attentive looking, of comparison and of argument, and we regard the revision as an improvement; we regard the preference of an inexperienced listener for a sentimental ballad over a Schubert song as a lapse of taste rather than as a difference of fact. Principle~\ref{p:super} is the near-universal view that two objects alike in every non-aesthetic respect, and properly perceived in the same category, cannot differ in beauty \citep{Sibley1959,Zangwill2001}. Principle~\ref{p:noncirc} is a general requirement on accounts of any concept.

\section{The argument}

\begin{lemma}\label{l:subj}
Subjectivism violates Principle~\ref{p:correct}.
\end{lemma}

\begin{proof}
On subjectivism, the content of $s$'s judgment that $x$ is beautiful is that $x$ is beautiful for $s$, which holds just in case $x$ gives $s$ aesthetic pleasure. A sincere judgment made by a judge to whom $x$ gives aesthetic pleasure is therefore true. And when $s$ judges $x$ beautiful and $t$ judges it not beautiful, $s$ says that $x$ is beautiful for $s$ and $t$ says that $x$ is not beautiful for $t$; both may be true together, so neither need be mistaken. Both clauses of Principle~\ref{p:correct} fail.
\end{proof}

Lemma~\ref{l:subj} is the familiar objection that subjectivism cannot make sense of aesthetic disagreement and of the improvement of taste. Hume put it by observing that whoever would equate Ogilby with Milton would be thought to defend an extravagance, although on the subjectivist view there is nothing to defend \citep{Hume1757}. The objection is not that subjectivism contradicts usage, which a revisionary theory may do, but that it removes the point of the practices of criticism and education in which aesthetic judgments are made.

\begin{lemma}\label{l:real}
Response-independent realism cannot account for Principle~\ref{p:answer}.
\end{lemma}

\begin{proof}
On response-independent realism, the possession of beauty does not consist in, and is not a priori connected to, any disposition to affect observers (Definition~\ref{d:real}). Then the coherence of an object that is beautiful and disposed to please no observer, however qualified, cannot be ruled out a priori: whatever connection holds between the property and our pleasure is a further fact, known if at all a posteriori, as the connection between a mineral's crystal structure and its appearance is known. Principle~\ref{p:answer} says that the connection between beauty and pleasure is a priori. The realist can therefore endorse it only by adding an a priori link to responses, and an account that adds such a link to the nature of the property is no longer response-independent.
\end{proof}

The realist may reply that the a priori link is a link between the \emph{concept} of beauty and pleasure, not between the property and pleasure: we pick out beauty by the pleasure it causes, as we might pick out a natural kind by its appearance, while its nature is independent of us. This reply concedes the point at issue for the concept, and with it the possibility of the realist's own scenario. If the concept fixes its reference by pleasure in qualified observers, then an object that pleases no qualified observer in any suitable condition does not fall under the concept, whatever intrinsic property it shares with objects that do. The realist property, if there is one, is then a property that beauty happens to coincide with, as the reflectance profile is a property that red happens to coincide with in the actual world. We return to this point under Objection~\ref{o:zang}.

\begin{proposition}\label{p:rdworks}
The response-dependent account, with aesthetic pleasure as the response, qualified observers as the observers and suitable conditions as the conditions (Definition~\ref{d:qual}), satisfies Principles~\ref{p:answer}--\ref{p:noncirc}.
\end{proposition}

\begin{proof}
The account says that it is a priori that $x$ is beautiful iff $x$ is disposed to give aesthetic pleasure to qualified observers in suitable conditions.

\emph{Answerability.} An object disposed to please no observer however qualified is not disposed to please qualified observers, and so is not beautiful; this is known a priori from the biconditional.

\emph{Correctability.} A judge whose own pleasure is the basis of her judgment may fail to be qualified, or may judge in unsuitable conditions. Her pleasure is then evidence about the object's disposition, but defeasible evidence, and the judgment may be false. When two judges disagree, they make incompatible claims about a single disposition of the object, and at least one of them is mistaken.

\emph{Supervenience.} Dispositions supervene on their categorical bases. Two objects alike in every non-aesthetic respect and properly perceived in the same category present the same stimulus to qualified observers in suitable conditions, and have the same disposition to please them.

\emph{Non-circularity.} Aesthetic pleasure is characterized by its object and its disinterestedness, and qualified observers and suitable conditions by the Humean conditions, none of which uses the concept of beauty.
\end{proof}

\begin{proposition}\label{p:exhaust}
Subjectivism, response-independent realism and response-dependence exhaust the accounts of beauty that satisfy Principle~\ref{p:super}.
\end{proposition}

\begin{proof}
Either the possession of beauty is a priori connected to dispositions to elicit responses in observers, or it is not. If it is not, the account is response-independent realism. If it is, either the relevant observers are each judge in turn, so that the property is relative to the judge, which is subjectivism; or they are observers of a kind specified independently of any particular judge. In the last case, Principle~\ref{p:noncirc} requires that the kind be specified without the concept of beauty, which yields a response-dependent property in the sense of Definition~\ref{d:rd}. An account on which the kind is specified using the concept of beauty violates Principle~\ref{p:noncirc} and is not a candidate.
\end{proof}

\begin{theorem}\label{t:main}
Beauty is a response-dependent property: an object is beautiful if and only if it is disposed to give aesthetic pleasure to qualified observers in suitable conditions.
\end{theorem}

\begin{proof}
By Proposition~\ref{p:exhaust}, any account satisfying the four principles is subjectivist, response-independent realist or response-dependent. By Lemma~\ref{l:subj}, it is not subjectivist; by Lemma~\ref{l:real}, it is not response-independent realist. It is therefore response-dependent, and by Proposition~\ref{p:rdworks} the account with Humean observers and conditions satisfies the principles.
\end{proof}

\begin{corollary}\label{c:beholder}
Beauty is not in the eye of the beholder: no individual's response either constitutes or guarantees an object's beauty.
\end{corollary}

\begin{proof}
An individual's response settles whether an object is beautiful only if the individual is qualified and in suitable conditions, and even then it is one instance of a disposition whose manifestation in other qualified observers in suitable conditions is equally relevant. No individual's response constitutes the property, and the response of an unqualified judge does not guarantee it.
\end{proof}

\begin{corollary}\label{c:relation}
Beauty is located in the relation between objects and a kind of observer. It is a property of the object, but a relational one: its instances are fixed by the object's non-aesthetic features together with the sensibility of qualified observers.
\end{corollary}

The position is the aesthetic analogue of colour relationalism, on which colours are real properties of objects that depend on perceivers without being illusory (MET-17). It agrees with the realist that judgments of beauty are true or false and are about objects; it agrees with the subjectivist that beauty would not exist in a world without sensibility. What it denies is that the sensibility in question is the beholder's.

\section{Objections and replies}

\begin{objection}[Euthyphro]\label{o:euth}
Qualified observers are pleased by beautiful objects because the objects are beautiful. Beauty is therefore prior to the response, and the response-dependent biconditional gets the order of explanation backwards.
\end{objection}

\begin{reply}
The objection assumes that qualified observers respond to beauty as such. On the account defended here, they respond to the non-aesthetic features of the object: to its proportions, colours, rhythms and their relations, discriminated with the delicacy that Definition~\ref{d:qual} requires. The explanation of a qualified observer's pleasure therefore runs through the object's non-aesthetic features and the observer's sensibility, and beauty is not one of its terms. In Wright's terms, the responses of qualified observers are \emph{extension-determining} for beauty, not merely extension-reflecting \citep{Wright1992}. The Euthyphro contrast is real for properties such as shape, where best judgments track a fact constituted independently of them; the claim of this manuscript is that beauty is not such a property, and the objection presents no consideration against that claim beyond the contrast itself.
\end{reply}

\begin{objection}[Circularity]\label{o:circ}
Hume's qualified critic is in practice identified by the soundness of her verdicts. Delicacy of taste is the capacity to perceive beauties others miss, and the account is circular.
\end{objection}

\begin{reply}
Delicacy can be characterized as the capacity to discriminate non-aesthetic features: to hear the intonation of a chord, to see the difference between two greens, to notice a recurring motif. Hume's own illustration, Sancho's kinsmen detecting the taste of leather and iron in the wine, is of this kind \citep{Hume1757}. Practice and comparison are matters of exposure and can be specified by history. Freedom from prejudice is the absence of interests in the object's author, owner or effects. That we often \emph{identify} good critics by their verdicts is an epistemic fact about how we find qualified observers, not a constitutive fact about what makes them qualified, and the account requires non-circularity only in the second.
\end{reply}

\begin{objection}[Divergence]\label{o:div}
Hume conceded that differences of temperament and age, and the manners and opinions of different periods, may produce variation in taste that no condition of qualification removes. If qualified observers diverge, the biconditional of Theorem~\ref{t:main} yields no verdict, and beauty is relative after all \citep{Goldman1995,Levinson2001}.
\end{objection}

\begin{reply}
Where qualified observers in suitable conditions diverge about an object, the account implies that the object neither determinately has nor determinately lacks the disposition to please them, and that the question whether it is beautiful has no determinate answer. This is a local indeterminacy, not a relativization of beauty to individuals. It leaves untouched every case in which qualified observers agree, which includes the cases on which the practice of criticism rests. Whether the area of divergence is large or small is an empirical question about the convergence of ideal critics, which AES-03 answers in favour of convergence. The account of the nature of beauty given here does not depend on that answer, although its practical significance does.
\end{reply}

\begin{objection}[Phenomenology]\label{o:zang}
Beauty is experienced as a feature of the object, not as a relation to us. The response-dependent account misdescribes what we experience and makes beauty less objective than it seems \citep{Zangwill2001}.
\end{objection}

\begin{reply}
The account says that beauty is a property of the object, and so agrees that it is experienced as such. That the property is relational is not something experience reveals or denies: colours are experienced as features of surfaces, and the claim that they are relational properties is a theory of what is experienced, not a report of how it seems (MET-17). As for objectivity, the account delivers everything Principles~\ref{p:correct} and \ref{p:super} require: judgments that are true or false, disagreements in which one party may be wrong, and dependence on the object's features. What it denies is that beauty could exist in a world in which no possible observer could take pleasure in anything, and that denial is part of the concept, as Principle~\ref{p:answer} records.
\end{reply}

\begin{objection}[Unobserved and changed sensibility]\label{o:unobs}
A flower in an unexplored valley is beautiful though no one sees it. And if human sensibility were altered so that qualified observers came to enjoy only discord, the flower would not thereby cease to be beautiful. The account seems to imply otherwise.
\end{objection}

\begin{reply}
The first case is handled by the dispositional form of the account: the flower is beautiful because it would please qualified observers if they perceived it in suitable conditions. The second is handled by indexing the account to actual sensibility, as is standard for response-dependent accounts of colour and value \citep{Johnston1989,Wright1992}: $x$ is beautiful iff $x$ is disposed to please observers qualified \emph{as we actually are}. Under the counterfactual alteration, the flower is still beautiful, because it is still disposed to please actual qualified observers; what changes is that the altered observers no longer track beauty.
\end{reply}

\section{What the result does not establish}

The argument settles the metaphysical form of beauty and not its extension. It does not show which objects are beautiful, nor that qualified observers converge on any particular range of cases; the second question is the topic of AES-03. It does not show that beauty is the same in art and in nature: the account is neutral on the knowledge that a qualified observer must have to perceive an object in its relevant category, and AES-11 argues that for natural objects this knowledge is scientific rather than art-historical. It does not extend automatically to every aesthetic property. Thick properties such as elegance or garishness may be response-dependent in the same way, but the argument here relies on the specific link between beauty and pleasure recorded in Principle~\ref{p:answer}, and a parallel argument for another property would need a parallel premise. Finally, the account does not imply that beauty is valuable only as a source of pleasure. That the property is defined by a disposition to please does not settle what makes the pleasure worth having.

\section{Conclusion}

Beauty is not in the eye of the beholder, if that means that each person's pleasure settles what is beautiful for that person, since this makes aesthetic error and disagreement impossible. Nor is beauty a property that objects might have though no observer could ever take pleasure in them, since the link between beauty and pleasure is a priori. Beauty is a response-dependent property: an object is beautiful just in case it is disposed to please qualified observers in suitable conditions, where qualification and suitability are specified, as Hume specified them, without appeal to beauty. Beauty is thus located in the relation between object and observer, and judgments of beauty are correct or incorrect according to whether that relation holds.

\begin{thebibliography}{99}\small
\bibitem[Goldman(1995)]{Goldman1995} Goldman, A.~H. (1995). \emph{Aesthetic Value}. Westview Press.
\bibitem[Hume(1757)]{Hume1757} Hume, D. (1757). Of the standard of taste. In \emph{Four Dissertations}. A. Millar.
\bibitem[Johnston(1989)]{Johnston1989} Johnston, M. (1989). Dispositional theories of value. \emph{Proceedings of the Aristotelian Society, Supplementary Volume}, 63, 139--174.
\bibitem[Kant(1790)]{Kant1790} Kant, I. (1790). \emph{Critique of the Power of Judgment}. Trans. P. Guyer and E. Matthews. Cambridge University Press, 2000.
\bibitem[Levinson(2001)]{Levinson2001} Levinson, J. (2001). Aesthetic properties, evaluative force, and differences of sensibility. In E. Brady and J. Levinson (Eds.), \emph{Aesthetic Concepts: Essays After Sibley}. Oxford University Press.
\bibitem[McDowell(1985)]{McDowell1985} McDowell, J. (1985). Values and secondary qualities. In T. Honderich (Ed.), \emph{Morality and Objectivity} (pp.~110--129). Routledge \& Kegan Paul.
\bibitem[Pettit(1983)]{Pettit1983} Pettit, P. (1983). The possibility of aesthetic realism. In E. Schaper (Ed.), \emph{Pleasure, Preference and Value: Studies in Philosophical Aesthetics}. Cambridge University Press.
\bibitem[Santayana(1896)]{Santayana1896} Santayana, G. (1896). \emph{The Sense of Beauty}. Charles Scribner's Sons.
\bibitem[Sibley(1959)]{Sibley1959} Sibley, F. (1959). Aesthetic concepts. \emph{The Philosophical Review}, 68(4), 421--450.
\bibitem[Wright(1992)]{Wright1992} Wright, C. (1992). \emph{Truth and Objectivity}. Harvard University Press.
\bibitem[Zangwill(2001)]{Zangwill2001} Zangwill, N. (2001). \emph{The Metaphysics of Beauty}. Cornell University Press.
\end{thebibliography}

\end{document}
